\documentclass[12pt, titlepage]{article}
\usepackage[margin=1in]{geometry}
\usepackage{graphicx}
\usepackage{float}

%\usepackage{hyperref}

\title{\Huge Feedback Controls Spring Final Design Report \\ \normalsize{EGEE 4810 - Senior Design}}
\author{by \\ \Large TEAM WINGFEATHER \\ Aydan Vissing \\ Brady Schick \\ Gabriel Southwell \\ Simon Ross}
\date{\footnotesize{Due} \\ \normalsize{April 29, 2026}}


%\setlength{\topmargin}{-.25in}
%\setlength{\oddsidemargin}{.25in}
%\setlength{\evensidemargin}{.25in}
%\setlength{\textwidth}{6.0in}
%\setlength{\textheight}{9.0in}
%\renewcommand{\baselinestretch}{1.4}
%\usepackage{indentfirst} %Maybe keep this maybe not

\begin{document}
\maketitle

%FIXME remove comments

%FIXME improve abstract
\begin{abstract}
This project report summarizes the progress to date on project Wingfeather, intended to reimagine the Feedback Controls Laboratory experience at Cedarville University. The main goals of this project are to make the lab experience more engaging and relevant for students, simpler to work with, more reliable than previous solutions, and more cost-effective. To achieve these ends, progress has been made towards a cheap drone platform that is capable of communicating fundamental feedback controls concepts along with a simple software user interface. Up to this point, hardware has been developed for the fabrication of a printed circuit board with the necessary components. This board has been acquired, and sensor communication code is in the integration phase. A mathematical model has been created to describe the drone dynamics, and will also be integrated soon. As the project continues to mature, the rest of the selected components will be integrated onto a custom 3D-printed frame and the testing phase will begin. Ultimately, this project is on track to be completed by the end of next semester. The current status shows promise to transform the way students view Feedback Controls lab, relating classwork to concepts they are familiar with and making lab time more enjoyable.


\end{abstract}

\newpage
\tableofcontents
%\listoffigures
%\listoftables

\newpage
\section{Definition Statement}

This project seeks to redefine the Cedarville Feedback Controls Laboratory experience by designing a quadcopter drone to use as an entertaining platform to teach classroom concepts. 

\section{Introduction \& Background}

The current lab experience for EGME-4660 uses the Feedback Analogue Unit 33-110 and the Feedback Mechanical Unit 33-100. Although the labs using the equipment do cover the desired content, they are often tedious and perhaps confusing to new users, especially users unfamiliar with Feedback Controls concepts. The user interface to conduct the experiments is a circuit board with nodes and banana plug cords. The user can connect these nodes to create feedback circuits whose effect is then shown using the mechanical unit. As mentioned, understanding what is happening between the units (and even in the electrical circuits) can be challenging for students. This causes the applications of the class material to  seem boring and irrelevant, especially for mechanical engineering students who may not have extensive electronics background. In addition, the equipment is old and fragile - adjusting the power supply or wiring the power incorrectly can result in damage to the equipment, often requiring parts to be replaced. Since the equipment is out of production and thus quite expensive, the current repair method has been de-soldering and soldering components by hand. We plan to improve the lab experience by replacing the current equipment with cheap and robust drones that can be reprogrammed for the purposes of the lab with a user-friendly interface. In this way, the instructor can choose the complexity of the tasks given to students for the lab.

\section{Objectives Revisited}
\subsection{Marketing Requirements}

The following requirements define the end user experience as desired by the customer:
  \begin{itemize}
    \item The final product must focus on autonomous actuation and control.
    \item The laboratory experience must include the ability to implement a PID control design.
    \item The final product must be able to physically and graphically display a step response.
    \item The students must have the opportunity to measure a Bode plot in real life.
    \item The laboratory experience must include the ability to see and correct the effects of physical and observable disturbances and noise.
    \item The final product should showcase the results of controllers physically in real life and not just with numbers and data.
    \item The final product must have the ability to log data to a neutral format independent of real time visualizations.
    \item The software solution must include a widget to generate understandable plots of the data.
    \item The final laboratory experience must include the capability to perform a tracking experiment.
    \item The laboratory experience must take relevant safety precautions.
    \item The final platform must be able to implement open loop control and log this response. 
  \end{itemize}
\subsection{Engineering Requirements}
\subsubsection{Hardware Requirements}
% improve
To fulfil the original vision of a simple to assemble, low cost, and safe to operate drone platform, we designed a drone platform from the ground up. 

The drone features a PCB that comes mostly assembled from the factory with only a few components that need to be soldered. To save on the cost, the pcb screws directly into a 3d printed frame that holds the motors and battery. To make the hardware safer to operate, the drone frame features ducted propellers with recessed blades. This reduces the chance of the blades catching student's hair or hands. 

\subsubsection{Software Requirements}

To give our team a good starting point on the drone, we based the drone software on an existing project called Esp-drone from Espressif. Esp-drone provides proven flight logic, sensor fusion, state estimation, and communication protocols. Building upon this basis, we were able to more easily reach our other software goals.

For the students to interact with the drones, we created cross-platoform application that is able to flash drone firmware, upload configurations, and plot flight data. This app can be extended to include any new features. 

\subsection{Constraints}
\begin{itemize}
  \item The laboratory experiments must be able to be conducted in ENS 145 at Cedarville University.
  \item The total cost of design, test, and evaluation must not exceed \$600.
  \item The final laboratory experience must not endanger the safety of the students.
  \item The final laboratory solution must not not destroy Cedarville property.
  \item The final laboratory experience must not use MATLAB.
\end{itemize}

\section{Engineering Design Revisited}
\subsection{Hardware}
\subsubsection{Frame Design}
When designing the frame, we chose to use Fusion360 for our editor rather than SolidWorks. This is because Fusion360 has a free version that is available consistently and since SolidWorks is less available to students. The frame itself was designed with the idea of being extremely light while still protecting the drone. The frame is sturdy enough to fly while maintaining the integrity of all components fastened to it. The current version is our best version since it accomplishes all the tasks that the previous frames couldn’t, such as keeping the motor mounts intact while it was being removed from the print bed and cleaned of support. The frame additionally has: a place to slot in the battery with Velcro, mounting places for the TOF sensors, a mindful orientation pattern for the PCB to frame, and blade guards to protect not just the blades, but also people around the drone. 
\subsubsection{PCB Design}
\subsection{Drone Software}
The backbone of the project is the esp-drone project from espressif. Esp-drone provides proven flight logic, sensor fustion, state estimation, and communication protocols. It created a great basis for our project that could be further extended anyware that we needed. 

\subsubsection{Added Components}
\begin{enumerate}
  \item Board
    \begin{itemize}
      \item Defines the pins the esp32 uses for interacting with the motors, sensors, and other peripherals.
      \item Contains public functions for getting and iterable array for time of flight xshut pins and motor pwm pins.
    \end{itemize}
  \item Buzzer
    \begin{itemize}
      \item Esp-drone used a piezo buzzer for alerting the users of various things. When selecting our hardware we chose to use an electromagnetic buzzer. While it is limited in the tones it can produce, it is much louder than a traditional piezo buzzer. 
      \item The new buzzer driver has one public function for starting a beep pattern. It takes in a enum and plays the corresponding pattern. There is also a public function for stopping the current pattern.
      \item It uses a low priority FreeRTOS task that sleeps until it gets a new play command. Once the pattern is complete it sleeps again. 
      \item A pattern definition includes sound on and off times deined in milliseconds. Some tones like low battery can be set to repeat indefinitely. 
    \end{itemize}
  \item File Manager
    \begin{itemize}
      \item The file manager provides helper functions for the serial file transfer system and the scripting components of the drone. These functions include holding a file object open and allowing writing to the file.
      \item Additionally, the file manager contains a function (that could be moved in the future) that sets the system PID parameters from the user-side config file.
      \item The file manager is currently designed to manage specifically the config and script files provided by the user (or the flasher app), but could be expanded to manage more files in the future.
    \end{itemize}
  \item Storage
    \begin{itemize}
      \item The storage component sets up the littleFS file system on the drone. It mounts and formats the partition, allowing for storage of user files and data. 
    \end{itemize}
  \item UART Listener
    \begin{itemize}
      \item This component is designed to wait for the user-provided files to be transferred over the UART bus. When the files are sent from the host device, it uses functions from the file manager component to write them to littleFS. The drone alerts the host device that it is ready, and the host begins the transfer.
      \item The UART listener uses a set baud rate and receives small chunks of each file, one after another, acknowledging in between.
      \item The files are written chunk-by-chunk into their respective littleFS locations. 
      \item If the user presses start before a file is received, the current flight files on the drone are used.
    \end{itemize}
  \item Multi-TOF Sensors
    \begin{itemize}
      \item A big problem with our multi-TOF sensor setup is that they all have the same i2c address. So to be able to use more than one at all, we need to readdress the sensors. 
      \item Every VL53L1 sensor has an Xshut pin with a weak pull up attatched to it. When xshut is low the device is not powered. Using this, only one tof can be active at a time for easy readdressing. 
      \item The driver library that came with esp-drone had a simple bug in the readdressing code where the software was sending a 7 bit address but was overwriting the lowest bit. This was fixed with a simple bit shift. 
      \item The initialization code was also changed to automatically readdress with an incrementing counter. That way every new sensor that gets initialized will get a fress address.
    \end{itemize}
  \item State Estimation with TOF Sensors
    \begin{itemize}
      \item By default, esp-drone uses a downward facing tof sensor for z estimation. We decided to use more tof range sensors to also estimate x and y position. 
      \item There are two sensors on the front and one on the side. The front sensors are spaced out with the idea being that the difference in ranges could be used for estimating yaw. That feature was not implemented in the end. 
      \item When adding one of the front sensors and the side sensor into the state estimation, we were able to get reasonably good state estimates. Occasionally the data gets noisy but that might be fixable.
    \end{itemize}

\end{enumerate}

\subsubsection{Interfacing with the Drone}
Crazyflie has a system for interfacing with crazyflie drones that carries over to our project. There is a python library called cflib which fulfils many of the wireless requirements. For it work with our drones however you \textbf{must} use the esp-drone fork of cflib instead of the official one.

Once connected to the drone's WiFi, you can run a python script to send commands and receive telemetry. This very useful for getting live data while testing.

\subsection{Hardware}
\begin{enumerate}

\item Custom PCB
\begin{itemize}
\item Successfully designed and ordered the first version of the PCB
\item PCB provides support to drone frame, but the frame is not comprised by it
\end{itemize}

\item Microcontroller
\begin{itemize}
\item Model selected: ESP32\_WROOM\_32N8 
\item Fast microcontroller with plenty of RAM
\item 2 cores for supporting sensor I/O, motors, and multiple PID loops
\item Has an integrated antenna for wireless data transfer - see software section below
\item Buttons included to reset microcontroller and start script
\end{itemize}

\item Inertial Measurement Unit (IMU)
\begin{itemize}
\item Model selected: MPU-6050
\item Accurately senses rotation in 3 axes for calculating error
\item Includes both MEMS gyroscope and accelerometer for 6 axis measurements
\item Currently successfully receiving IMU data both wired and wirelessly.
\end{itemize}

\item Buzzer
\begin{itemize}
\item Model selected: FUET-8540-3.6V
\item Added buzzer to design for auditory feedback
\end{itemize}

\item LEDs
\begin{itemize}
\item Added 3 LED array for visual feedback
\end{itemize}

\item Connector
\begin{itemize}
\item Decided to use USB-C for connection to lab computers
\end{itemize}

\item Power Regulator
\begin{itemize}
\item Model selected: TPS63070RNMR
\item Decided to use buck/boost converter instead of a standard power regulator
\item Input range of 2-16V
\item Output voltage is 3.3V
\end{itemize}

\item Sensors
\begin{itemize}
\item Model selected: VL53L1
\item Mounting these sensors on the main PCB proved to be expensive and complicated, so we decided to use smaller boards to be soldered on to the main board
\item Range of 4 m, meaning we are able to see both the roof and floor at the same time. This will help with decision logic if the drone flies over an object
\item Decided to mount sensors on the top and bottom of the drone and 3 on the sides. This allows for accurate altitude measurements and ranging from walls. Additionally, by placing 2 sensors on one side of the drone, a heading may be calculated
\item Currently successfully receiving sensor data both wired and wirelessly.
\end{itemize}

\item Motors and Motor Controllers
\begin{itemize}
\item 720 coreless DC motors selected
\item Since the motors are brushed and only need to turn in one direction, MOSFETs are used instead of a full motor controller IC
\end{itemize}

\item Special Circuitry
\begin{itemize}
\item Diode circuit added to allow powering of the circuit by a battery, connected computer, or even both (without damage)
\item ESD and reverse current protection circuits included
\item Auto reset logic circuit included so the microcontroller enters the bootloader properly when a computer is connected
\end{itemize}

\item Battery
\begin{itemize}
\item 2 battery connector options included on the PCB - user may select which connector they are using and solder it on
\item The power converter accepts a large range of battery voltages, allowing for the use of 1, 2, or 3 cell LiPo batteries
\item Battery choices are flexible (helpful, considering needs may change). This way we can increase the voltage to the motors if needed.
\item Planning to need some sort of charging array for the students to each receive multiple batteries per lab session
\end{itemize}

\item Frame
\begin{itemize}
\item Cheap, light, strong 3D printed frame created that successfully holds the motors, battery, and printed circuit board
\item Adjustments being made to the print integrity and sensor placement on the frame
\end{itemize}

\end{enumerate}
\subsection{Controls}
\begin{enumerate}
		\item Non-Linear Model
		\begin{itemize}
			\item Fully defines how the drone will be controlled in software using PID controllers.
			\item PID controllers were used for their simplicity and their widespread use in controls fields
		\end{itemize}
		\item Linear Model
		\begin{itemize}
			\item Allows for mapping of poles of the control system
			\item This helps the tuning process to avoid producing an unstable system
		\end{itemize}
		\item Simulation
		\begin{itemize}
			\item Quad-copter dynamics and controls simulated using both MATLAB and python
			\item Allows for verification of quad-copter dynamics and controls
			\item Can act as a comparison of the real quad-copter to an ideal quad-copter
		\end{itemize}
		\item Crazyflie Controller
		\begin{itemize}
			\item Specifies the control system in the drone software
			\item Has reprogrammable values within the controller for tuning
			\item Includes outLimit and iLimit to improve drone stability
		\end{itemize}
		\item PID Tuning
		\begin{itemize}
			\item Developed process for fully tuning drone
			\item Has several parameters that can be logged and plotted for observability during the process
		\end{itemize}
\end{enumerate}

\section{Project Summary}

\subsection{Final Results}

\subsection{Challenges Overcome}


As of now, we have a fully designed PCB, an almost fully designed drone frame, and the majority of the parts needed to assemble the quadcopter. A lot of the work has also been done in setting up the drone firmware and designing the flight controls. Moving forward, we plan to assemble the drone and finish up the firmware and controls. Then, we will be able to move on to controller tuning, lab design, and setting up the desktop application.

\subsection{Future Potential}
\begin{enumerate}
\item Simulation
\begin{itemize}
\item Adding nested controls to simulation to match Crazyflie's controls
\item Fully converting MATLAB simulation into python
\item Implementing python simulation as a prediction model in the app
\end{itemize}
\end{enumerate}

\subsection{Budget}
The budget below describes exactly how we used the funds alotted to us and stayed under the \$600 maximum spending. 

\begin{table}[H]
\centering
\caption{Unit Budget for Feedback Controls Lab Prototype}
\label{tab:budget}
\begin{tabular}{||c|c|c|c||}
\hline 
\hline
\textbf{Component/Expense} & \textbf{Price} & \textbf{Comments} \\ 
\hline 
Prototype Microcontroller & \$0 & Dr. Kohl gave us one to borrow \\
\hline 
Time of Flight Sensors & \$45 & Original 2(\$15) + an additional 5* (\$30) \\
\hline 
Motors & \$10 & 4 motors \\ 
\hline 
Body & \$0 & 3D printed \\ 
\hline 
Props & \$17 & Two seperate orders, testing and application \\ 
\hline 
Battery & \$0 & Used batteries on hand \\ 
\hline 
Bolt and Wrench Sets & \$76 & Needed for the Gimble setup \\ 
\hline 
PCB(s) & \$67 & 2 boards populated boards, 3 unpopulated \\ 
\hline 
\textbf{Total Cost} & \textbf{\$215} &  &  \\ 
\hline 
\hline
\end{tabular} 
\end{table}


\end{document}
